\chapter{Ablation}

\section*{Definition and Overview}
Ablation is a group of medical techniques that deliberately destroy or remove unwanted tissue using heat (radiofrequency or laser), extreme cold (cryoablation), microwaves, ultrasound, or chemicals. Because the energy can usually be delivered through thin catheters, needles, or probes rather than an open incision, ablation is typically minimally invasive, offering a shorter recovery than conventional surgery. The technique is applied across many specialties, including cardiac electrophysiology for heart rhythm disorders, interventional radiology and surgical oncology for tumours, gynaecology for heavy menstrual bleeding, and vascular surgery for varicose veins. The specific approach, risks, and expected outcomes depend closely on the tissue being treated.

\section*{Epidemiology}
Catheter ablation has become a standard treatment for many arrhythmias. It is increasingly used for atrial fibrillation, the most common significant rhythm disorder, which affects roughly 1 in 30 adults and becomes more prevalent with advancing age. Endometrial ablation is a well-established alternative to hysterectomy for heavy menstrual bleeding, a problem affecting about a quarter of women of reproductive age. Percutaneous tumour ablation is widely used for primary and secondary cancers of the liver, kidney, and lung, particularly in people who are not candidates for surgical resection. Endovenous ablation is now a first-line treatment for symptomatic varicose veins.

\section*{Aetiology and Pathophysiology}
Ablation is used to treat a range of conditions in which localised destruction of tissue produces benefit.

\begin{itemize}
    \item \textbf{Cardiac arrhythmias:} abnormal electrical circuits or accessory pathways in the heart that cause palpitations, including supraventricular tachycardia, atrial flutter, and atrial fibrillation
    \item \textbf{Tumours:} primary or metastatic cancers of the liver, kidney, lung, or bone that are unsuitable for surgical removal
    \item \textbf{Heavy menstrual bleeding:} destruction of the endometrium in women who have completed their family
    \item \textbf{Varicose veins:} closure of incompetent veins by heat delivered from within the vein
    \item \textbf{Other applications:} treatment of Barrett's oesophagus, small thyroid nodules, some benign bone tumours, and selected pain or sinus conditions
\end{itemize}

Whatever the energy source, the principle is the same: controlled, localised tissue damage causes the abnormal tissue to scar, cease functioning, shrink, or die.

\section*{Clinical Features}
Symptoms before ablation reflect the underlying condition, such as palpitations, breathlessness, and fatigue from an arrhythmia, heavy and prolonged menstrual bleeding, or aching and swelling from varicose veins. After the procedure, certain features are expected during recovery.

\begin{itemize}
    \item Mild pain, bruising, or swelling at the access site is common and settles over days
    \item Fatigue for a few days is normal after most ablative procedures
    \item Palpitations may persist for weeks after cardiac ablation while the heart rhythm settles
    \item Light vaginal bleeding may continue for a few weeks after endometrial ablation
    \item The treated vein becomes firm and gradually fades and shrinks after vein ablation
\end{itemize}

\section*{Differential Diagnosis}
The alternatives to ablation are considered when planning treatment for each condition.

\begin{itemize}
    \item Medical therapy to control symptoms rather than an ablative procedure
    \item Conventional surgery, such as tumour resection or hysterectomy
    \item Active surveillance, for tumours that are small or slow-growing
    \item Different ablative technologies, such as radiofrequency, cryotherapy, or laser
    \item Other minimally invasive options, such as embolisation of tumours or uterine artery embolisation for heavy bleeding
\end{itemize}

\section*{When to Seek Medical Care}
After ablation, the treatment team should be contacted promptly if any of the following develop.

\begin{itemize}
    \item Fever, chills, or rigors
    \item Severe or worsening pain at the treatment site
    \item Heavy bleeding, or discharge that looks infected
    \item Redness, swelling, or heat spreading from the site
    \item After cardiac ablation: chest pain, severe breathlessness, or a return of worrying palpitations
    \item After vein ablation: calf pain or swelling, or breathing difficulty, which may indicate a blood clot
\end{itemize}

\textbf{Call 000 (or local emergency services) for:} collapse, severe chest pain, difficulty breathing, or sudden severe bleeding.

\section*{Diagnosis and Investigations}
Each application of ablation requires condition-specific work-up before the procedure.

\begin{itemize}
    \item \textbf{Cardiac ablation:} an electrocardiogram and electrophysiology study to map the abnormal electrical pathway before ablation
    \item \textbf{Tumour ablation:} CT, MRI, or ultrasound to localise and size the tumour and to plan image-guided needle placement
    \item \textbf{Endometrial ablation:} pelvic ultrasound and, in appropriate cases, sampling of the endometrium to exclude cancer beforehand
    \item \textbf{Vein ablation:} duplex ultrasound of the leg veins to map the refluxing vein and confirm the deep veins are healthy
    \item \textbf{Pre-procedure blood tests} and a medication review, with particular attention to anticoagulant and antiplatelet drugs
    \item \textbf{Follow-up imaging or review} to confirm the intended effect and detect recurrence
\end{itemize}

\section*{Management}
Management comprises the choice of ablative modality, the technical performance of the procedure, and appropriate aftercare.

\begin{itemize}
    \item \textbf{Catheter ablation (heart):} thin catheters are advanced from the groin to the heart under sedation or anaesthesia, and the abnormal tissue is heated or frozen
    \item \textbf{Tumour ablation:} a needle or probe is guided into the tumour through the skin under imaging control, and heat, cold, or microwaves destroy the cancer
    \item \textbf{Endometrial ablation:} a small probe passed through the cervix destroys the lining of the womb
    \item \textbf{Endovenous ablation:} a laser or radiofrequency catheter heats the varicose vein from within, sealing it shut
    \item \textbf{Pain relief and rest} in the period after the procedure
    \item \textbf{Follow-up} to assess success and to manage any recurrence or complication
\end{itemize}

\section*{Complications}
\begin{itemize}
    \item Bleeding, bruising, or infection at the access site
    \item Pain and swelling in the treated area
    \item Damage to structures adjacent to the treatment site, which varies with location
    \item After cardiac ablation: transient arrhythmias, vascular injury, and rarely stroke or pericardial effusion
    \item After endometrial ablation: fluid absorption, uterine perforation, and, if pregnancy occurs subsequently, a risk of serious placental complications
    \item After tumour ablation: bleeding, injury to nearby organs, or pneumothorax when the lung is treated
    \item Incomplete treatment with recurrence of the condition, sometimes requiring a repeat procedure
\end{itemize}

\section*{Prognosis and Outcomes}
Outcomes depend strongly on the condition treated. Catheter ablation cures the majority of supraventricular tachycardias and controls atrial fibrillation effectively in many patients, although some require a repeat procedure. Endometrial ablation reduces bleeding in most women, though a minority later need further intervention or hysterectomy. Percutaneous ablation achieves good local control in appropriately selected tumours, and endovenous ablation permanently closes the treated vein in the great majority of cases. Experienced teams and careful patient selection are the most important determinants of success.

\section*{Prevention and Self-Care}
\begin{itemize}
    \item Discuss the options, success rates, and risks with a specialist before deciding on ablation
    \item Tell the treatment team about all medications, especially blood thinners
    \item Follow instructions on fasting and on stopping or continuing medicines before the procedure
    \item Rest and avoid heavy activity for the recommended period afterwards
    \item Report any new or unusual symptoms promptly
    \item After endometrial ablation, use contraception until pregnancy is certainly not possible, because pregnancy after the procedure is dangerous
    \item Maintain a healthy lifestyle to reduce the chance of the underlying condition, such as atrial fibrillation, recurring
\end{itemize}

\section*{Sources and Further Reading}
The following organisations provide reliable information on ablation in its various applications.

\begin{itemize}
    \item NHS. \textit{Radiofrequency ablation: overview of the procedure.} https://www.nhs.uk/conditions/radiofrequency-ablation/
    \item Mayo Clinic. \textit{Cardiac catheter ablation: overview and risks.} https://www.mayoclinic.org/tests-procedures/cardiac-ablation/
    \item MedlinePlus. \textit{Ablation therapy health information.} https://medlineplus.gov/ablationtherapy.html
    \item MSD Manual. \textit{Overview of arrhythmias and treatment options.} https://www.msdmanuals.com/
    \item British Heart Foundation. \textit{Catheter ablation for heart rhythm disorders.} https://www.bhf.org.uk/
    \item NICE. \textit{Guidance on interventional procedures including ablation.} https://www.nice.org.uk/guidance/conditions-and-diseases/
\end{itemize}
